\chapter{Compactness and balanced subspaces}

\begin{lemma}[Bounded lattice bases]
\label{lem:bounded-lattice-bases}\notready
Fix $n\geq1$, $0<c\leq C$ and $\varepsilon>0$. Full lattices
$\Lambda\subset\mathbb R^n$ with covolume in $[c,C]$ and
$\|x\|\geq\varepsilon$ for $0\ne x\in\Lambda$ admit bases whose
vectors have norms bounded by a constant depending only on
$n,c,C,\varepsilon$.
\end{lemma}
\begin{proof}
Let $v_n$ be the volume of the unit ball and put
$R=3(C/v_n)^{1/n}$. A ball of radius $R/2$ has volume greater than
$\covol\Lambda$. It contains two distinct points differing by a lattice
vector: otherwise its translates into a half-open fundamental
parallelepiped would be disjoint almost everywhere, contradicting the
volume inequality. More explicitly, integrate over that parallelepiped
the integer-valued sum of indicators of all translates of the ball;
its integral exceeds the parallelepiped's volume, so its value exceeds
one at some point. Their difference is a nonzero lattice vector of norm
less than $R$.

Choose a shortest nonzero vector $u$, of length $\ell\in[\varepsilon,R]$.
It is primitive, for divisibility by an integer of absolute value greater
than one would give a shorter vector. Extend it to an integral basis.
Orthogonal projection to $u^\perp$ gives a full lattice $\Lambda'$
with covolume $\covol\Lambda/\ell\in[c/R,C/\varepsilon]$.
Every nonzero projected vector has length at least
$\sqrt3\ell/2$: lift it and subtract an integer multiple of $u$ so
the longitudinal component has absolute value at most $\ell/2$;
the resulting nonzero lattice vector has length at least $\ell$.
Induction in dimension therefore supplies a bounded basis of
$\Lambda'$. Lift each basis vector and again adjust its longitudinal
component into $[-\ell/2,\ell/2]$. Together with $u$ these lifts form
a basis of $\Lambda$, and their norms are bounded by the square root
of the squared projected bound plus $R^2/4$. In dimension one the
single generator has norm equal to covolume, at most $C$, starting
the induction.
\end{proof}

\begin{lemma}[Compact lattice frames]
\label{lem:mahler-frames}\uses{lem:bounded-lattice-bases}\notready
Under the hypotheses of the preceding lemma, there is a compact set
of nonsingular real $n\times n$ matrices containing a basis matrix of
every such lattice.
\end{lemma}
\begin{proof}
Let $M$ be the common bound for basis vectors. The set of matrices whose
columns have norm at most $M$ and whose absolute determinant lies in
$[c,C]$ is closed and bounded in $\mathbb R^{n^2}$, hence compact.
The positive lower determinant bound makes every matrix nonsingular,
and the preceding lemma puts a basis matrix of each lattice in this set.
\end{proof}

\begin{lemma}[Arithmetic compactness]
\label{lem:arithmetic-compactness}
\uses{def:hermitian-degree,lem:restriction-covolume,lem:mahler-frames}
\source{bost-specht-manuscript:page-0011}\notready
For fixed $K$ and $a\geq1$, every sequence of normalized semistable
rank-$a$ Hermitian $\Ocal$-bundles has a subsequence which, after
$\Ocal$-module isomorphisms, has one fixed underlying module and
convergent metrics. The limit is normalized and semistable.
\end{lemma}
\begin{proof}
For $0\ne x\in A$, the free line $\Ocal x$ has degree
$-D^{-1}\sum_v w_v\log\|x_v\|_v$. Semistability and normalization
therefore give $\prod_v\|x_v\|_v^{w_v}\geq1$. The arithmetic--geometric
mean inequality, applied to $D$ positive numbers with each complex
norm repeated twice, gives $\sum_vw_v\|x_v\|_v^2\geq D$.
The underlying Euclidean lattice has systole at least $\sqrt D$ and
covolume $\Delta_K^{a/2}$. Compact lattice frames give integral bases
whose Gram matrices converge to a positive-definite matrix $G$.

For $\alpha\in\Ocal$, multiplication by $\alpha$ has operator norm at
most $\max_v|\alpha_v|$. Convergence of the positive-definite Gram
matrices bounds their condition numbers and hence bounds all entries
of these integral multiplication matrices. Choose a fixed integral
basis $\alpha_1,\ldots,\alpha_D$ of $\Ocal$. There are only finitely
many possibilities for the associated $D$ matrices, so pass to a
subsequence on which all are constant. These matrices fix the entire
$\Ocal$ action, and identify every module in the subsequence with one
fixed module. The limit is still a positive-definite metric.

This metric comes from the required archimedean metrics. Indeed, the
orthogonal idempotents of $K\otimes_\mathbb Q\mathbb R$ are real
linear combinations of the fixed multiplication matrices, and are
self-adjoint in every metric. These are closed matrix equations, so
they persist in the limit. The components remain orthogonal. On each
complex component multiplication by $i$ is orthogonal, another closed
equation, and thus the restriction is a Hermitian metric. Determinants
are continuous, so the covolume formula gives normalization. For each
fixed submodule its determinant norm varies continuously, and its
degree remains at most zero in the limit. Thus all semistability
inequalities persist.
\end{proof}

\begin{lemma}[The maximum degree and openness of strict stability]
\label{lem:degree-maximum-continuity}
\uses{lem:degree-finiteness,lem:restriction-covolume}
\source{bost-specht-manuscript:page-0011,bost-specht-manuscript:page-0012}
\notready
For every Hermitian bundle $A$, the number
\[
 \mathcal D(A)=\max\{\degn(B):B\subseteq A\text{ saturated}\},
 \qquad\degn(0)=0,
\]
exists. It is continuous in the metrics. The strictly stable normalized
metrics on a fixed module form an open subset of the normalized metrics.
\end{lemma}
\begin{proof}
Degree finiteness shows that the degrees at least zero come from a
finite set of submodules, including zero, so their maximum exists and
dominates every other degree. If two sets of norms differ by factors
between $e^{-t}$ and $e^t$, a rank-$b$ determinant norm changes by a
factor between $e^{-bt}$ and $e^{bt}$. Each normalized degree therefore
changes by at most $bt$, and $\mathcal D$ changes by at most $at$.
Convergence of positive-definite metrics gives such bounds with $t\to0$,
proving continuity, also along the fixed-module subsequences above.

At a normalized strictly stable metric every nonzero proper saturated
submodule has negative degree. Degree finiteness applied with lower
bound $-1$ gives a finite set of exceptional submodules; those outside
the set have degree less than $-1$. Choose $\delta>0$ such that every
exceptional nonzero proper submodule has degree at most $-2\delta$,
and $2\delta\leq1$. The same bound then holds for all nonzero proper
submodules. A metric change with $at<\delta$ keeps all their degrees
negative. Restricting to normalized metrics proves openness.
\end{proof}

\begin{lemma}[Intrinsic partial traces]
\label{lem:balanced-partial-traces}
\source{bost-specht-manuscript:page-0012}
\lean{BostConjecture.Balanced.partialTraceRight, BostConjecture.Balanced.partialTraceLeft, BostConjecture.Balanced.trace_partialTraceRight_comp, BostConjecture.Balanced.trace_partialTraceLeft_comp, BostConjecture.Balanced.trace_comp_tensorSum_of_partialTrace_eq_smul}\leanok
Let $U,W$ be finite-dimensional vector spaces over a field $k$.
Under the canonical isomorphism
$\operatorname{End}(U\otimes W)=\operatorname{End}(U)\otimes
\operatorname{End}(W)$, define $\Tr_W=1\otimes\Tr$ and
$\Tr_U=\Tr\otimes1$. For $P\in\operatorname{End}(U\otimes W)$,
$H\in\operatorname{End}(U)$ and $J\in\operatorname{End}(W)$,
\[
 \Tr((\Tr_W P)H)=\Tr(P(H\otimes1)),\qquad
 \Tr((\Tr_U P)J)=\Tr(P(1\otimes J)).
\]
Consequently, if $\Tr_W P=a1$ and $\Tr_U P=b1$, then
$\Tr(P(H\otimes1+1\otimes J))=a\Tr H+b\Tr J$.
\end{lemma}
\begin{proof}
The tensor-endomorphism isomorphism sends $A\otimes B$ to the map
$u\otimes w\mapsto Au\otimes Bw$. Such maps span all endomorphisms:
in bases of $U,W$, their matrix units are exactly the matrix units on
the tensor-product basis. Both sides of each claimed identity are linear
in $P$, so it suffices to take $P=A\otimes B$. Then
$P(H\otimes1)=AH\otimes B$ and
$\Tr(AH\otimes B)=\Tr(AH)\Tr(B)$, by summing diagonal entries in
the tensor-product basis. This equals $\Tr((\Tr(B)A)H)$.
The second identity follows from
$P(1\otimes J)=A\otimes BJ$ and the same diagonal calculation.
Adding the identities and substituting the scalar partial traces gives
the last formula.
\end{proof}

\begin{lemma}[Invariant detection of orbit closure]
\label{lem:invariant-orbit-criterion}\notready
For a product $G$ of complex special linear groups acting algebraically and linearly on
a finite-dimensional complex vector space $X$, a nonzero $x\in X$
has $0\notin\overline{Gx}^{\mathrm{Zar}}$ if and only if some
positive-degree homogeneous $G$-invariant polynomial is nonzero at $x$.
The Euclidean and Zariski closures of $Gx$ coincide.
\end{lemma}

\begin{lemma}[Balanced projections give geometric semistability]
\label{lem:balanced-plucker}
\uses{def:geometric-semistability,lem:invariant-orbit-criterion,lem:balanced-partial-traces}
\source{bost-specht-manuscript:page-0012}\notready
Let $S\subset U\otimes W$ be a nonzero subspace of complex Hermitian
spaces of dimensions $m,n$, and put $s=\dim S$. If its orthogonal
projection $P$ obeys
$\Tr_W P=(s/m)\Id_U$ and $\Tr_U P=(s/n)\Id_W$, then its Pl\"ucker
line is geometrically semistable for $\SL(U)\times\SL(W)$.
\end{lemma}
\begin{proof}
Take an orthonormal basis $v_1,\ldots,v_s$ of $S$ and
$z=v_1\wedge\cdots\wedge v_s$, so $\|z\|=1$.
For traceless Hermitian $H,J$, the derivative at zero of
$\|\bigwedge^s(e^{tH}\otimes e^{tJ})z\|^2$ is
\[
 2\Tr(P(H\otimes\Id+\Id\otimes J))=0.
\]
To see the derivative, differentiate the determinant of the Gram
matrix of the transformed $v_i$ at the identity Gram matrix; its
derivative is the sum of the diagonal derivatives, giving the trace
displayed. The induced infinitesimal operator on the exterior power
is Hermitian, so diagonalization writes the squared norm as
$\sum_r c_r e^{2t\lambda_r}$ with $c_r\geq0$ and real $\lambda_r$.
Its second derivative is nonnegative, and its zero derivative at
$t=0$ makes $t=0$ a global minimum.

For $g\in\SL(U)$, polar decomposition writes $g=u e^H$, where $u$
is unitary and $H$ is traceless Hermitian: the positive factor has
determinant $|\det g|=1$. The same applies to $W$. Unitary factors
preserve exterior norms, so every vector in the orbit of $z$ has norm
at least one. Thus zero is outside its Euclidean closure and, by the
orbit criterion, outside its Zariski closure. The same criterion gives
a nonvanishing homogeneous invariant, as required.
\end{proof}
